\documentclass[11pt]{article}
\usepackage[margin=1.2in]{geometry}
\usepackage{amsmath,amssymb,amsthm}

\theoremstyle{plain}
\newtheorem{theorem}{Theorem}[section]
\newtheorem{lemma}[theorem]{Lemma}
\newtheorem{corollary}[theorem]{Corollary}
\theoremstyle{remark}
\newtheorem{remark}[theorem]{Remark}

\newcommand{\R}{\mathbb{R}}
\newcommand{\Dg}{\Delta g}

\setlength{\parindent}{0pt}
\setlength{\parskip}{0.85\baselineskip}

\title{\large Additive Residual Coordinates and Gradient Updates}
\author{}
\date{}

\begin{document}
\maketitle

% =============================================================================
\section{Additive Residual Coordinates}
% =============================================================================

Residual pruning is efficient only when the score used by the decoder admits a
local update.  If deleting an element changes the relevant residual statistic by
a fixed local correction, then the score coordinate must preserve this additive
structure.  The following elementary facts formalize this obstruction.

\begin{theorem}[Exact residual updates force affine coordinates]
  \label{thm:affine}
  Let $\mathcal{X}$ be a state space and let
  $\varphi:\mathcal{X}\to\R^p$ be a score encoding.  Suppose that for every
  elementary residual transformation $T_\alpha$ there exists a fixed vector
  $c_\alpha\in\R^p$ such that
  \[
    \varphi(T_\alpha x)=\varphi(x)+c_\alpha
  \]
  for all admissible states $x$.  Then, for every state obtained from a
  reference state $x_0$ by a sequence $T_{\alpha_r}\cdots T_{\alpha_1}$,
  \[
    \varphi(T_{\alpha_r}\cdots T_{\alpha_1}x_0)
    =
    \varphi(x_0)+\sum_{q=1}^r c_{\alpha_q}.
  \]
  In particular, if the residual state is represented by coordinates
  $d\in\mathbb{Z}^m$ and elementary transformations induce fixed increments
  $\Delta_\alpha$, then $\varphi$ is affine on the generated lattice:
  \[
    \varphi(d)=Cd+b.
  \]
\end{theorem}

\begin{proof}
  The first identity follows by repeated substitution of the assumed update
  rule.  In coordinate form, the relation
  $\varphi(d+\Delta_\alpha)-\varphi(d)=c_\alpha$ says that all discrete first
  differences along the generated directions are constant.  Hence $\varphi$ is
  affine on each connected lattice component generated by these directions.
\end{proof}

\begin{theorem}[Bounded nonlinear outputs cannot carry nontrivial exact updates]
  \label{thm:nonlinear}
  Let $\sigma:\R\to\R$ be continuous and bounded, and suppose a coordinate of
  the score has the form $\varphi_j(x)=\sigma(z_j(x))$.  If some residual
  transformation has a nonzero additive correction in coordinate $j$ and can be
  applied along arbitrarily long trajectories, then $\varphi_j$ cannot satisfy
  the exact additive update rule along all such trajectories.
\end{theorem}

\begin{proof}
  If the update rule held for a transformation with correction
  $(c_\alpha)_j\ne 0$, then along $x_r=T_\alpha^r x_0$ we would have
  \[
    \varphi_j(x_r)=\varphi_j(x_0)+r(c_\alpha)_j,
  \]
  which is unbounded as $r\to\infty$.  This contradicts boundedness of
  $\sigma(z_j(x_r))$.
\end{proof}

\begin{theorem}[Approximate residual scores induce ordering errors at scale]
  \label{thm:approx_flip}
  Let a decoder rank elements by scalar residual scores $s_i^{(n)}$ on an
  instance of size $n$. Suppose an approximating model returns
  $\hat{s}_i^{(n)}=s_i^{(n)}+\varepsilon_i^{(n)}$. Let
  \[
    \gamma_n := \min_{i\ne j}|s_i^{(n)}-s_j^{(n)}|
  \]
  be the minimum nonzero score gap. If $\gamma_n\to 0$ along a sequence of
  instance sizes and the approximation error is not $o(\gamma_n)$, then no such
  approximation can guarantee preservation of the deletion order for all large
  instances. In particular, there exist arbitrarily large instances on which a
  pairwise ordering is inverted.
\end{theorem}

\begin{proof}
  Correct deletion order requires every pairwise score gap to be preserved. Thus,
  for all $i,j$ with $s_i^{(n)}<s_j^{(n)}$, one must have
  \[
    \hat{s}_i^{(n)} < \hat{s}_j^{(n)}.
  \]
  Equivalently,
  \[
    \varepsilon_i^{(n)}-\varepsilon_j^{(n)}
    <
    s_j^{(n)}-s_i^{(n)}.
  \]
  The smallest possible right-hand side is $\gamma_n$. Hence a uniform guarantee
  requires approximation errors below the minimum-margin scale. If the error is
  not $o(\gamma_n)$, then for arbitrarily large $n$ the error can be of the same
  order as the closest score gap, and an adverse error direction reverses that
  pair. Therefore an approximating decoder cannot guarantee the exact deletion
  order at scale.
\end{proof}

\begin{corollary}[Vanishing margins make fixed-accuracy approximations insufficient]
  \label{cor:vanishing_margin}
  In residual pruning problems where relevant score gaps shrink with $n$, a
  model whose score error is controlled only to fixed accuracy cannot guarantee
  the exact deletion trajectory for all large $n$. Exact additive coordinates
  avoid this failure mode by updating the residual score algebraically rather
  than repeatedly approximating it.
\end{corollary}

\begin{theorem}[A local ordering error induces trajectory divergence]
  \label{thm:cascade}
  Consider an iterative decoder that selects the element with minimum score
  criterion. If, at some step, approximate scores invert the ordering of two
  candidate elements, then the approximate and exact decoders delete different
  elements and subsequently evolve on different residual states. Later score
  discrepancies are therefore governed by structural divergence of the states,
  not by the original local approximation error.
\end{theorem}

\begin{proof}
  Once the ordering is inverted, the two decoders choose different elements. The
  residual states after the deletion are consequently distinct. Every later
  score is then evaluated on a different residual structure, so the next
  discrepancy is a difference between two scoring problems, not a perturbation
  of the same score vector. Such discrepancies can compound across the
  trajectory.
\end{proof}

% =============================================================================
\section{Gradient Coordinates}
% =============================================================================

The obstruction above does not imply that useful additive coordinates are
unavailable.  For local objectives, the objective gradient itself carries a
residual correction signal.  We write the objective in the form
\[
  \mathcal{L}(p;G)
  =
  \sum_i \phi_i(p_i)
  +
  \sum_{(i,j)\in F(G)} \psi_{ij}(p_i,p_j,A_{ij}),
\]
where $p_i$ is the scalar value attached to element $i$, $F(G)$ is the set of
local factors in the current residual instance, and $A_{ij}$ denotes the local
relation between $i$ and $j$.  Define
\[
  g_i(G)=\frac{\partial \mathcal{L}(p;G)}{\partial p_i}.
\]

\begin{lemma}[Local deletion gives a local gradient correction]
  \label{lem:local_gradient_correction}
  If the residual transition deletes element $v$, removing the local factor
  between $i$ and $v$, then for every surviving element $i$,
  \[
    g_i(G\setminus v)-g_i(G)
    =
    -\frac{\partial \psi_{iv}(p_i,p_v,A_{iv})}{\partial p_i},
  \]
  with the convention that the right-hand side is zero when no such factor is
  present.
\end{lemma}

\begin{proof}
  By differentiating the decomposed objective,
  \[
    g_i(G)=\phi_i'(p_i)+
    \sum_{j:(i,j)\in F(G)}
    \frac{\partial \psi_{ij}(p_i,p_j,A_{ij})}{\partial p_i}.
  \]
  Deleting $v$ leaves all terms in the sum unchanged except the term involving
  $v$.  The stated correction follows immediately.
\end{proof}

\begin{theorem}[Gradient injection exposes an additive pre-activation coordinate]
  \label{thm:additive_coordinate}
  Let a hidden unit receive an ordinary feature vector $x$ and a scalar gradient
  coordinate $g$ through an affine pre-activation
  \[
    z(g)=w^\top[x;g]+b.
  \]
  If a residual transition changes $g$ by a computable amount $\delta$, then
  \[
    z(g+\delta)=z(g)+w_g\delta,
  \]
  where $w_g$ is the coefficient multiplying $g$.  Thus the pre-activation is
  an exact additive coordinate for the gradient correction.
\end{theorem}

\begin{proof}
  Since only the scalar input $g$ changes,
  \[
    z(g+\delta)=w^\top[x;g+\delta]+b
    =w^\top[x;g]+b+w_g\delta
    =z(g)+w_g\delta.
  \]
\end{proof}

\begin{remark}[Representability and training]
  The previous theorem is a representability statement.  It shows that gradient
  injection provides a coordinate in which local residual corrections can be
  expressed additively.  It does not imply that standard training necessarily
  identifies this coordinate or uses it consistently over long deletion
  trajectories.  A model may learn a useful static score while failing to place
  the relevant signal in an updateable coordinate.
\end{remark}

% =============================================================================
\section{Gradient Updater}
% =============================================================================

The Gradient Updater (GU) makes the gradient signal available during decoding,
rather than relying on the learned static score to carry it implicitly.  At each
residual step, the updater receives the current score, the removed vertex, the
local relation to that vertex, and a live gradient-derived coordinate.  This
turns the local correction identified in Lemma~\ref{lem:local_gradient_correction}
into an operational update signal.

\begin{theorem}[GU has sufficient local information for residual correction]
  \label{thm:gu_information}
  For a locally decomposable objective, the effect of deleting $v$ on the
  gradient coordinate of a surviving element $i$ is determined by the local
  factor involving $(i,v)$.  Consequently, an updater supplied with the live
  gradient state and the local relation $A_{iv}$ has the information required to
  represent the residual correction.
\end{theorem}

\begin{proof}
  Lemma~\ref{lem:local_gradient_correction} gives the correction as the removed
  local derivative
  \[
    -\frac{\partial \psi_{iv}(p_i,p_v,A_{iv})}{\partial p_i}.
  \]
  This quantity depends only on the local relation between $i$ and $v$ and on
  the current values appearing in that factor.  Providing the updater with the
  live gradient signal and local relation therefore supplies the information
  needed to model the correction.
\end{proof}

\begin{remark}
  GU should be viewed as an optimization/identifiability device.  Gradient
  injection makes an additive coordinate representable; GU keeps the corresponding
  residual signal explicit during decoding, reducing the burden on SGD to encode
  it perfectly inside a nonlinear static score.
\end{remark}

% =============================================================================
\section{Fast Gradient Maintenance}
% =============================================================================

A direct implementation that recomputes the dense gradient after every deletion
costs $O(n^2)$ per step and $O(n^3)$ over a full trajectory.  For local
objectives, this recomputation is unnecessary: deleting a vertex removes only
that vertex's contribution to each surviving gradient coordinate.

\subsection*{General form}

For each surviving element $i$, the gradient can be written as
\[
  g_i
  =
  \sum_j \psi^{\nabla}_{ij}(p_i,p_j,A_{ij})
  +
  \text{unary or simple global terms}.
\]
After deleting $v$, all summands with $j\ne v$ remain unchanged.  Hence
\[
  g_i' = g_i - \psi^{\nabla}_{iv}(p_i,p_v,A_{iv}),
\]
possibly together with a simple update to any global scalar term.  The update
is therefore a vector correction over surviving vertices, not a new dense
matrix-vector computation.

\subsection*{Clique objective}

For the clique loss
\[
  \mathcal{L}_{\mathrm{clique}}(p;A,k)
  =
  -\frac{p^\top A p}{k(k-1)}
  +
  \frac{p^\top B p}{k(n-k)}
  +
  10\left(\frac{\sum_i p_i-k}{n}\right)^2,
\]
where $B$ is the complement adjacency matrix with zero diagonal, the gradient is
\[
  g_i
  =
  -\frac{2(Ap)_i}{k(k-1)}
  +
  \frac{2(Bp)_i}{k(n-k)}
  +
  \frac{20(\sum_j p_j-k)}{n^2}.
\]
It is therefore enough to cache
\[
  q^A=Ap,
  \qquad
  q^B=Bp,
  \qquad
  S=\sum_j p_j.
\]
Deleting vertex $v$ with value $p_v$ gives
\[
  q^A \leftarrow q^A-A_{:v}p_v,
  \qquad
  q^B \leftarrow q^B-B_{:v}p_v,
  \qquad
  S \leftarrow S-p_v.
\]
The corrected gradient is then read from the same formula using the updated
cached quantities.

\begin{theorem}[Cached gradient maintenance yields quadratic decoding]
  \label{thm:cached_gradient_complexity}
  On a dense graph, suppose one GNN forward pass costs $O(n^2)$. The initial
  gradient cache $(Ap,Bp,\sum_i p_i)$ can be built in $O(n^2)$ by two
  matrix-vector products. If each deletion updates this cache by subtracting the
  deleted vertex's column contribution, then each residual step costs $O(n)$ and
  the full decoding trajectory costs $O(n^2)$.
\end{theorem}

\begin{proof}
  The initial forward pass and the initial products $Ap$ and $Bp$ cost
  $O(n^2)$.  At a deletion step, updating $Ap$ and $Bp$ requires subtracting two
  length-$n$ vectors, $A_{:v}p_v$ and $B_{:v}p_v$, and updating $S$ requires one
  scalar subtraction.  Thus each deletion costs $O(n)$.  Since there are at most
  $n$ deletions, the total update cost is $O(n^2)$, and the initial setup is of
  the same order.
\end{proof}

\begin{remark}[Cached deletion gradient]
  The cache maintains the gradient induced by deleting vertices from the current
  residual objective under the cached probability values.  If the learned updater
  also changes all probabilities after each deletion, recomputing the exact
  fresh dense gradient of those changed probabilities would again require a
  dense matrix-vector product.  The cached gradient should therefore be understood
  as a deletion-updated residual signal used by GU, not as a full recomputation
  after arbitrary global score drift.
\end{remark}


The algebraic obstruction explains why a nonlinear static score need not be an
exact residual coordinate.  The gradient observation shows that the local
correction signal is nevertheless available: deleting an element removes a local
contribution from the objective gradient.  Gradient injection makes this signal
representable, while GU keeps it explicit during decoding.  With cached gradient
maintenance, the live gradient signal can be updated in $O(n)$ per deletion,
giving $O(n^2)$ total decoding after the initial GNN pass.

\end{document}
